\subsection{K-Combinações}

Na lib padrão do c++, tem a função next\_permutation:

\inccode{comb/go_through_perms.cpp}

Dados inteiros positivos $N$ e $K$, determinar os subconjuntos de $\{1, \dots, N\}$ de tamanho exatamente $K$.

\prob{1}{kcombs1} Determine a próxima K-combinação (considerando ordem lexicográfica) considerando a primeira como sendo $1, 2, \dots, K$

\textbf{Solução:} Em $O(K)$

\inccode{comb/next_comb.cpp}

\prob{2}{kcombs2} Gere todas as K-combinações de modo que combinações adjacentes difiram por no máximo um elemento

\textbf{Solução 1:} Em $O(2^n)$. Usa o fato de que se você gerar os $2^n$ graycodes de $N$ bits e pegar apenas os de $K$ bits ativos, então a diferença entre caras consecutivos vai ser que de um pro outro um bit apaga e um bit acende.

\inccode{comb/adjacent_kcombs.cpp}

\textbf{Solução 2:} Em $O\left( N \cdot {N \choose K} \right)$, mas com constante alta

\inccode{comb/adjacent_kcombs_high_const.cpp}
